\unit{Coherent, number and thermal light}
\unitref{Weeks 1--2; photon statistics tutorial; Assignment 1}
\subsection{What does classical-looking light mean?}
A coherent state satisfies \(\hat a\ket\alpha=\alpha\ket\alpha\)
for a complex amplitude \(\alpha=|\alpha|e^{i\theta}\).
Unlike a number state, it is a \emph{superposition}
of different photon numbers:
\[
\boxed{\ket\alpha=e^{-|\alpha|^2/2}
\sum_{n=0}^\infty\frac{\alpha^n}{\sqrt{n!}}\ket n.}
\]
The probability of \(n\) photons is obtained by
squaring the corresponding amplitude:
\[
P(n)=e^{-|\alpha|^2}\frac{|\alpha|^{2n}}{n!}.
\]
This is Poissonian: \(\langle n\rangle=|\alpha|^2\),
\(\operatorname{Var}(n)=|\alpha|^2\).
A coherent state has a mean oscillating electric
field because \(\langle a\rangle=\alpha\), with
\(\langle E\rangle\propto2|\alpha|\cos(kz-\omega t+\theta)\).

\subsection{Displacement: making a coherent state}
The unitary displacement operator
\[
D(\alpha)=e^{\alpha a^\dagger-\alpha^*a}
\]
acts on vacuum to give \(\ket\alpha=D(\alpha)\ket0\).
A useful calculation employs the BCH identity:
\[
D^\dagger aD=a+[a,\alpha a^\dagger-\alpha^*a]
+\tfrac12[[a,\alpha a^\dagger-\alpha^*a],\ldots]
=\boxed{a+\alpha}.
\]
The first commutator equals \(\alpha\), a scalar;
all higher nested commutators vanish. Therefore
\(aD\ket0=D(a+\alpha)\ket0=\alpha D\ket0\),
which proves the eigenstate property.

\subsection{A density operator describes mixtures}
A pure state \(\ket\psi\) has \(\rho=\ket\psi\bra\psi\).
When preparation involves several alternative states,
\[
\boxed{\rho=\sum_jp_j\ket{\psi_j}\bra{\psi_j},
\quad \langle O\rangle=\operatorname{Tr}(\rho O).}
\]
A mixture differs from a coherent superposition.
For example,
\((\ket0+\ket1)/\sqrt2\) has off-diagonal
density-matrix terms, whereas a 50--50 mixture
of \(\ket0,\ket1\) does not. Both have the
same probabilities in a photon-number measurement,
but different phase-sensitive expectations.
Purity is measured by \(\operatorname{Tr}\rho^2\):
it equals 1 for a pure state and is smaller for a mixture.

\subsection{Thermal light and bunching}
A single thermal optical mode has
\[
\rho_{\rm th}=\sum_{n=0}^\infty
\frac{\bar n^n}{(1+\bar n)^{n+1}}\ket n\bra n,
\quad
\bar n=\frac{1}{e^{\hbar\omega/(k_BT)}-1}.
\]
The variance is
\(\operatorname{Var}(n)=\bar n+\bar n^2\).
Compare with coherent light, where the variance
equals the mean, and a fixed-\(n\) Fock state,
where the variance vanishes.
The Mandel parameter
\[
Q_M=\frac{\operatorname{Var}(n)-\bar n}{\bar n}
\]
describes super-Poissonian (\(Q_M>0\)),
Poissonian (\(Q_M=0\)) and sub-Poissonian
(\(Q_M<0\)) fluctuations.

\begin{worked}{Three states with mean photon number 4}
A coherent state with \(|\alpha|^2=4\)
has variance \(4\). A thermal mode with
\(\bar n=4\) has variance \(4+16=20\).
The Fock state \(\ket4\) has variance \(0\).
Equal mean energy does \emph{not} mean equal noise.
\end{worked}

\subsection{Second-order coherence counts pairs}
The zero-delay Glauber correlation is
\[
\boxed{g^{(2)}(0)=
\frac{\langle a^{\dagger2}a^2\rangle}{\bar n^2}
=\frac{\langle n(n-1)\rangle}{\bar n^2}}.
\]
Why \(n(n-1)\)? It measures the number of
ordered pairs that can be formed from \(n\)
photons, which is relevant for two detections.
A one-photon state has \(g^{(2)}=0\):
it cannot cause two ideal simultaneous counts.
A coherent state has \(g^{(2)}=1\);
a thermal mode has \(g^{(2)}=2\).

\subsection{Entangled pairs can look thermal locally}
A two-mode squeezed state has the form
\[
\ket\Psi=\sqrt{1-\lambda^2}
\sum_n\lambda^n\ket n_s\ket n_i.
\]
It has perfectly correlated photon numbers.
Tracing over the unobserved idler gives
\[
\rho_s=(1-\lambda^2)\sum_n\lambda^{2n}\ket n\bra n,
\quad \bar n=\frac{\lambda^2}{1-\lambda^2}.
\]
This is \emph{thermal locally} despite the joint state
being pure. The signal--idler cross-correlation is
\[
g_{si}^{(2)}(0)=
\frac{\langle n_sn_i\rangle}{\bar n^2}
=2+\frac1{\bar n}.
\]
At low gain, detection of one idler photon
heralds one signal photon.

\begin{checkpoint}
A photon-number histogram alone cannot reveal
all coherences or entanglement. Ask whether
the experiment measures number, field phase,
or correlations between modes.
\end{checkpoint}
